\section{Discussion}
\label{sec:discussion}

\subsection{Interpretation}

SENTINEL demonstrates that multimodal environmental monitoring through
AI-driven fusion provides earlier and more specific pollution detection
than any single modality. The fusion system (AUROC = 0.939) outperforms
the best individual encoder (AquaSSM, AUROC = 0.920), confirming that
cross-modal information is complementary rather than redundant. The key
insight is that different modalities have fundamentally different response
latencies: behavioral signals respond within minutes, sensor anomalies
within hours, microbial community shifts within days, and satellite-visible
changes within weeks. Fusion leverages this temporal hierarchy for both
rapid initial detection and sustained diagnostic specificity.

ToxiGene represents, to our knowledge, the first supervised method for
multi-label toxicity classification from RNA-seq. Its P-NET biological
hierarchy provides interpretability absent from black-box alternatives:
parameter attribution traces predictions to specific genes, pathways, and
biological processes, enabling mechanistic insight into detected
contamination. The 31.3$\times$ temporal persistence ratio and FPR = 0.000
on clean reference sites confirm operational reliability.

\subsection{Limitations}

Several limitations should be noted:

\begin{itemize}
    \item \textbf{Data availability}: Truly multimodal data (all five
    modalities at the same location-time) is rare. SENTINEL-DB contains
    390M+ records but co-located five-modality samples are sparse. Most
    training relies on single-modality or modality-pair data.

    \item \textbf{Geographic bias}: 95\% of training data comes from
    the United States (USGS, EPA) and Europe (GRQA). Performance in
    tropical, arid, or data-sparse regions is expected to degrade.

    \item \textbf{MicroBiomeNet scope}: While F1 = 0.913 on real EMP 16S
    data validates the Aitchison-geometry design, the 8-class source
    classification task is simpler than fine-grained pollution source
    attribution. More targeted aquatic microbiome datasets with
    contaminant-specific labels would better evaluate diagnostic utility.

    \item \textbf{HydroViT WQ regression}: While MAE pretraining converges,
    the 16-parameter water quality regression head requires spatiotemporally
    co-registered satellite--in situ pairs, which are sparse in SENTINEL-DB.

    \item \textbf{Computational cost}: The full 189.5M-parameter system
    requires GPU acceleration for training; inference is more modest
    ($\sim$50ms per forward pass on A100).

    \item \textbf{BioMotion chronic endpoints}: BioMotion achieves 0\%
    detection on PFAS and pharmaceutical exposures, which produce chronic
    rather than acute behavioral effects. These endpoints fall outside the
    training distribution of acute ECOTOX assays.

    \item \textbf{EPA regulatory correlation}: Correlation between
    SENTINEL alert scores and EPA violation counts is non-significant
    ($\rho = 0.112$, $p = 0.858$), likely because EPA violations reflect
    compliance sampling frequency rather than true contamination prevalence.
\end{itemize}

\subsection{Broader Impact}

SENTINEL has potential for significant positive environmental impact by
enabling earlier detection of water pollution events, reducing exposure
of both ecosystems and human populations. The public platform component
democratizes access to water quality intelligence. However, we note that
AI-based monitoring should complement, not replace, direct chemical
analysis and regulatory oversight.

\subsection{Future Work}

A preliminary discovery scan across 18 US water bodies flagged 8 of 9
known-impaired sites as ALERT, providing early evidence of deployment
readiness. Future directions include: (1) extending to additional water
body types (groundwater, coastal, deep ocean), (2) active learning for
optimal sampling strategies, (3) integration with real-time alert systems,
(4) deployment in resource-limited settings through model distillation,
and (5) extending BioMotion training to chronic exposure assays for PFAS
and pharmaceutical detection.
